\documentclass[12pt]{article}
\usepackage{amsmath}
\usepackage{enumitem}
\usepackage[margin=1in]{geometry}
\usepackage{siunitx}
\usepackage{booktabs}

\title{Worked Solutions: Proportion Extension Problems}
\date{}

\begin{document}

\maketitle

\section*{Direct Proportion Past}

\begin{enumerate}
    \item \textbf{A bag of 5 apples costs £1.20. What is the cost of 40\% more apples? How many apples is that?}
    
    \textbf{Solution:}
    \begin{align*}
        \text{Cost per apple} &= \frac{£1.20}{5} = £0.24 \\
        40\% \text{ more apples} &= 5 + (40\% \times 5) = 5 + 2 = 7 \text{ apples} \\
        \text{Cost of 7 apples} &= 7 \times £0.24 = £1.68
    \end{align*}
    So 7 apples cost £1.68.
    
    \item \textbf{If 10\% of a number is 15, what is 40\% of the number?}
    
    \textbf{Solution:}
    \begin{align*}
        10\% \text{ of number} &= 15 \\
        1\% \text{ of number} &= 15 \div 10 = 1.5 \\
        40\% \text{ of number} &= 40 \times 1.5 = 60
    \end{align*}
    Alternatively: Since 40\% is 4 times 10\%, \(40\% = 4 \times 15 = 60\).
    
    \item \textbf{A shop offers a 15\% discount for buying 3 of the same item. If one book normally costs £8, how much would 3 books cost with the discount?}
    
    \textbf{Solution:}
    \begin{align*}
        \text{Normal cost for 3 books} &= 3 \times £8 = £24 \\
        \text{Discount amount} &= 15\% \times £24 = 0.15 \times 24 = £3.60 \\
        \text{Final cost} &= £24 - £3.60 = £20.40
    \end{align*}
    
    \item \textbf{The cost of fuel is directly proportional to the number of litres purchased. If 40 litres cost £68.00, and the price then increases by 5\%, what will 25 litres cost at the new price?}
    
    \textbf{Solution:}
    \begin{align*}
        \text{Original price per litre} &= \frac{£68.00}{40} = £1.70 \\
        \text{New price per litre} &= £1.70 \times 1.05 = £1.785 \\
        \text{Cost for 25 litres} &= 25 \times £1.785 = £44.625 \approx £44.63
    \end{align*}
    
    \item \textbf{A coffee shop's revenue is directly proportional to the number of customers. On Monday, they had 120 customers and made £420. On Tuesday, they increased their prices by 8\%. How many customers did they have on Tuesday if their revenue was £453.60?}
    
    \textbf{Solution:}
    \begin{align*}
        \text{Monday: revenue per customer} &= \frac{£420}{120} = £3.50 \\
        \text{Tuesday: new price per customer} &= £3.50 \times 1.08 = £3.78 \\
        \text{Number of customers Tuesday} &= \frac{£453.60}{£3.78} = 120
    \end{align*}
    They had 120 customers on Tuesday as well.
\end{enumerate}

\section*{Direct Proportion Future}

\begin{enumerate}[resume]
    \item \textbf{The perimeter of a square is directly proportional to the length of its side. If a square with a side of 4 cm has a perimeter of 16 cm, what is the perimeter of a square with a side of 11 cm?}
    
    \textbf{Solution:}
    \begin{align*}
        P &= 4 \times \text{side} \\
        \text{For side 11 cm: } P &= 4 \times 11 = 44 \text{ cm}
    \end{align*}
    
    \item \textbf{The length and width of a rectangle are in a direct proportion. A rectangle has a length of 5 cm and a perimeter of 18 cm. What is the perimeter of a rectangle with a length of 15 cm?}
    
    \textbf{Solution:}
    \begin{align*}
        \text{Original: } P &= 2(l + w) = 18 \\
        2(5 + w) &= 18 \\
        5 + w &= 9 \\
        w &= 4 \text{ cm}
    \end{align*}
    Since length and width are in direct proportion, when length triples from 5 cm to 15 cm, width also triples:
    \[
    \text{New width} = 3 \times 4 = 12 \text{ cm}
    \]
    \[
    \text{New perimeter} = 2(15 + 12) = 2 \times 27 = 54 \text{ cm}
    \]
    
    \item \textbf{The area of a square is directly proportional to the square of its side length. If a square with a side of 3 cm has an area of 9 cm², what is the area of a square with a side of 7 cm?}
    
    \textbf{Solution:}
    \begin{align*}
        A &= s^2 \\
        \text{For } s = 7 \text{ cm: } A &= 7^2 = 49 \text{ cm}^2
    \end{align*}
    
    \item \textbf{Two similar triangles are shown below. The perimeter of the smaller triangle is 21 cm. The side lengths of the larger triangle are 2.5 times those of the smaller one. What is the perimeter of the larger triangle?}
    
    \textbf{Solution:}
    For similar shapes, perimeter scales by the same factor as side lengths:
    \[
    \text{Larger perimeter} = 2.5 \times 21 = 52.5 \text{ cm}
    \]
    
    \item \textbf{A rectangular garden is 10 m long and 6 m wide. A new, similarly shaped garden is to be built. The area of the new garden needs to be 50\% larger than the original.}
    \begin{enumerate}[label=(\alph*)]
        \item \textbf{What is the area of the new garden?}
        \begin{align*}
            \text{Original area} &= 10 \times 6 = 60 \text{ m}^2 \\
            \text{New area} &= 60 \times 1.5 = 90 \text{ m}^2
        \end{align*}
        
        \item \textbf{What is the perimeter of the new garden?}
        
        Since the gardens are similar, the linear scale factor \(k\) satisfies:
        \[
        k^2 = \frac{\text{New area}}{\text{Original area}} = \frac{90}{60} = 1.5
        \]
        \[
        k = \sqrt{1.5} = \sqrt{\frac{3}{2}} = \frac{\sqrt{6}}{2}
        \]
        Original perimeter \(= 2(10 + 6) = 32 \text{ m}\)
        \[
        \text{New perimeter} = k \times 32 = \frac{\sqrt{6}}{2} \times 32 = 16\sqrt{6} \text{ m} \approx 39.19 \text{ m}
        \]
    \end{enumerate}
\end{enumerate}

\section*{Rates of Change Past}

\begin{enumerate}[resume]
    \item \textbf{A plant was 20 cm tall. Over the next week, it grew by 5\%. At what speed did it grow, in cm per day?}
    
    \textbf{Solution:}
    \begin{align*}
        \text{Growth amount} &= 5\% \times 20 = 0.05 \times 20 = 1 \text{ cm} \\
        \text{Speed} &= \frac{1 \text{ cm}}{7 \text{ days}} \approx 0.1429 \text{ cm/day}
    \end{align*}
    
    \item \textbf{A bathtub fills at a rate of 8 L/min. The plug is slightly loose, and water drains out at a rate that is 2\% of the fill rate. If the tub starts empty, how much water is in it after 10 minutes?}
    
    \textbf{Solution:}
    \begin{align*}
        \text{Drain rate} &= 2\% \times 8 = 0.16 \text{ L/min} \\
        \text{Net fill rate} &= 8 - 0.16 = 7.84 \text{ L/min} \\
        \text{Water after 10 min} &= 7.84 \times 10 = 78.4 \text{ L}
    \end{align*}
    
    \item \textbf{A car's value was £15,000. It depreciates at a rate of 12\% per year. A motorbike's value was £8,000 and it depreciates at 8\% per year. After 3 years, which vehicle is worth more?}
    
    \textbf{Solution:}
    Using depreciation formula: Value after \(n\) years = Initial \(\times (1 - \text{rate})^n\)
    \begin{align*}
        \text{Car after 3 years} &= 15000 \times (1 - 0.12)^3 = 15000 \times 0.88^3 \\
        &= 15000 \times 0.681472 = £10222.08 \\
        \text{Bike after 3 years} &= 8000 \times (1 - 0.08)^3 = 8000 \times 0.92^3 \\
        &= 8000 \times 0.778688 = £6229.50
    \end{align*}
    The car is worth more (£10222.08 vs £6229.50).
    
    \item \textbf{A social media post is shared rapidly. On the first day, it gets 200 shares. The number of shares increases by 15\% each day after that.}
    \begin{enumerate}
        \item \textbf{How many shares does it get on the third day?}
        \begin{align*}
            \text{Day 1:} & 200 \\
            \text{Day 2:} & 200 \times 1.15 = 230 \\
            \text{Day 3:} & 230 \times 1.15 = 264.5 \approx 265 \text{ shares}
        \end{align*}
        
        \item \textbf{What is the total number of shares at the end of the third day?}
        \[
        200 + 230 + 265 = 695 \text{ shares}
        \]
    \end{enumerate}
    
    \item \textbf{A factory's energy costs are £10,000 per month. The manager invests in new equipment which is 20\% more efficient. However, due to increased production, the factory now runs for 15\% more hours each month, up from 200 hours. Calculate the new energy cost per hour, and per month.}
    
    \textbf{Solution:}
    \begin{align*}
        \text{Original cost per hour} &= \frac{£10000}{200} = £50 \\
        \text{New efficiency: cost per hour} &= £50 \times (1 - 0.20) = £50 \times 0.8 = £40 \\
        \text{New running hours} &= 200 \times 1.15 = 230 \text{ hours} \\
        \text{New monthly cost} &= £40 \times 230 = £9200
    \end{align*}
\end{enumerate}

\section*{Rates of Change Future}

\begin{enumerate}[resume]
    \item \textbf{A square has a side length that is increasing at a rate of 2 cm/s. What is the rate of change of its perimeter when the side length is 5 cm? What about when the side length is 10 cm? Does the rate depend on the side length?}
    
    \textbf{Solution:}
    \begin{align*}
        P &= 4s \\
        \frac{dP}{dt} &= 4 \frac{ds}{dt} = 4 \times 2 = 8 \text{ cm/s}
    \end{align*}
    This rate is constant and does not depend on the side length. So at both \(s = 5\) cm and \(s = 10\) cm, \(\frac{dP}{dt} = 8\) cm/s.
    
    \item \textbf{A rectangular car park is 60 m long and 40 m wide. It is to be resurfaced. The cost of resurfacing is £15 per square metre. Calculate the total cost.}
    
    \textbf{Solution:}
    \begin{align*}
        \text{Area} &= 60 \times 40 = 2400 \text{ m}^2 \\
        \text{Cost} &= 2400 \times 15 = £36000
    \end{align*}
    
    \item \textbf{A circular oil spill is spreading. It starts with a radius of 1 m. The radius of the spill is increasing at a constant rate of 0.5 m/min.}
    \begin{enumerate}
        \item \textbf{Calculate the rate at which the circumference is increasing when the radius is 4 m.}
        \begin{align*}
            C &= 2\pi r = \pi D \\
            \frac{dC}{dt} &= 2\pi \frac{dr}{dt} = 2\pi \times 0.5 = \pi \text{ m/min}
        \end{align*}
        This rate is constant, so at \(r = 4\) m, \(\frac{dC}{dt} = \pi \approx 3.14\) m/min.
        
        \item \textbf{Calculate the rate at which the circumference is increasing when the radius is 8 m.}
        \[
        \frac{dC}{dt} = \pi \text{ m/min} \text{ (same as above)}
        \]
        
        \item \textbf{Find the area of the spill after 10 minutes.}
        \begin{align*}
            \text{Radius after 10 min} &= 1 + 0.5 \times 10 = 6 \text{ m} \\
            A &= \pi r^2 = \pi \times 6^2 = 36\pi \approx 113.10 \text{ m}^2
        \end{align*}
    \end{enumerate}
    
    \item \textbf{A gardener is planting a rectangular flower bed. The length is increasing at 1 m/h and the width at 0.5 m/h. At a specific moment, the length is 6 m and the width is 4 m.}
    \begin{enumerate}
        \item \textbf{At what rate is the perimeter increasing at this moment?}
        \begin{align*}
            P &= 2(l + w) \\
            \frac{dP}{dt} &= 2\left(\frac{dl}{dt} + \frac{dw}{dt}\right) = 2(1 + 0.5) = 3 \text{ m/h}
        \end{align*}
        
        \item \textbf{How much will the area have increased by 30 minutes later? Do you think it grew at a constant rate?}
        
        Initial area \(A_0 = 6 \times 4 = 24 \text{ m}^2\).
        
        After 0.5 hours:
        \begin{align*}
            l &= 6 + 1 \times 0.5 = 6.5 \text{ m} \\
            w &= 4 + 0.5 \times 0.5 = 4.25 \text{ m} \\
            A_{\text{new}} &= 6.5 \times 4.25 = 27.625 \text{ m}^2 \\
            \text{Increase} &= 27.625 - 24 = 3.625 \text{ m}^2
        \end{align*}
        The area does NOT grow at a constant rate because:
        \[
        \frac{dA}{dt} = l \frac{dw}{dt} + w \frac{dl}{dt}
        \]
        which depends on both \(l\) and \(w\), which are changing.
    \end{enumerate}
    
    \item \textbf{A decorative path is being built around a rectangular pond. The pond is 10 m long and 6 m wide. The path has a uniform width of \(x\) metres all the way around.}
    \begin{enumerate}
        \item \textbf{Show that the area of the path, \(A\), is given by \(A = 4x^2 + 32x\).}
        
        \textbf{Solution:}
        \begin{align*}
            \text{Outer rectangle dimensions:} & \quad \text{length} = 10 + 2x, \quad \text{width} = 6 + 2x \\
            \text{Area of outer rectangle} &= (10 + 2x)(6 + 2x) = 60 + 20x + 12x + 4x^2 = 60 + 32x + 4x^2 \\
            \text{Area of pond} &= 10 \times 6 = 60 \\
            \text{Area of path} &= (60 + 32x + 4x^2) - 60 = 4x^2 + 32x
        \end{align*}
        
        \item \textbf{If the path is being paved at a rate of 0.8 m²/min, how wide is the path after 20 minutes?}
        
        \textbf{Solution:}
        \begin{align*}